% Fungal BGC Atlas manuscript, prepared in the style of the journal Database (Oxford Academic).
% NOTE: this uses a clean standard article-class preamble, not OUP's proprietary
% submission class file (not safely reproducible offline) -- swap in the official
% oup-authoring-template class at submission time if the journal requires it.
\documentclass[11pt]{article}

\usepackage[margin=2.5cm]{geometry}
\usepackage[utf8]{inputenc}
\usepackage[T1]{fontenc}
\usepackage{times}
\usepackage{graphicx}
\usepackage{booktabs}
\usepackage{caption}
\usepackage{amsmath}
\usepackage[colorlinks=true, linkcolor=blue, citecolor=blue, urlcolor=blue]{hyperref}
\usepackage[numbers,sort&compress]{natbib}
\usepackage{authblk}
\usepackage{lineno}
\usepackage{setspace}
\usepackage{xcolor}

\newcommand{\todo}[1]{\textcolor{red}{\textbf{[TODO: #1]}}}

\title{\textbf{The Fungal BGC Atlas: a curated, evidence-linked knowledge base\\
of 609 literature- and MIBiG-derived fungal biosynthetic\\
gene cluster dossiers}}

\author[1]{Richard Barker\thanks{Correspondence: \texttt{dr.richard.barker@gmail.com}}}
\affil[1]{University of Wisconsin--Madison, Madison, WI, USA}

\date{}

\begin{document}
\doublespacing
\linenumbers
\maketitle
\begin{center}
\textit{Article type: Database. Prepared in the format of the journal \emph{Database} (Oxford Academic).}
\end{center}

\begin{abstract}
\noindent
Fungal biosynthetic gene clusters (BGCs) underlie a chemically and biologically important class of
secondary metabolites, yet the primary evidence for any given cluster --- which genes were actually
shown to be required, which metabolites were directly observed, under what conditions, and how
confidently --- is scattered across decades of literature and only partially captured by curated
sequence repositories such as MIBiG. We present the \textbf{Fungal BGC Atlas}, a curated,
evidence-linked knowledge base of 609 fungal BGC dossiers, each anchored to a MIBiG accession and
built by large-language-model-assisted extraction from primary literature followed by full manual
review of every dossier. Beyond the cluster-level record, the Atlas resolves each dossier into 1572
genes, 807 metabolites, 915 experiments, 936 measurements, 1534 claims and 1264 curated
subject--predicate--object relationships, all traceable to 2575 source records through 47\,768
explicit evidence links. Structural and functional context is captured along 27 independent facet
axes (biosynthesis class, chemistry, ecology, genome context and activation conditions), and 995
conflicts or unresolved uncertainties are preserved as first-class records rather than collapsed
into a forced consensus. We package the Atlas as a relational SQLite database, a knowledge graph,
and a self-contained interactive web dashboard (Overview, Taxonomy, Chemistry, Facet Explorer,
Knowledge Graph, Evidence, Conflicts, a searchable Dossier Browser with cross-table detail views, and
a live Data Dictionary), deployed on GitHub Pages and archived on Zenodo with an open licence
(CC-BY-4.0). The Atlas turns 609 individually FAIR-adjacent literature records into a single
queryable, comparable, and citable resource for fungal natural-product genome mining.\\[0.5em]
\textbf{Database URL:} \url{https://dr-richard-barker.github.io/fungal-bgc-atlas/}
\end{abstract}

\textbf{Keywords:} biosynthetic gene clusters; fungal natural products; MIBiG; knowledge graph;
FAIR data; curated database; secondary metabolism; evidence provenance


\newpage
\section{Introduction}

Fungi are prolific producers of secondary metabolites -- polyketides, non-ribosomal peptides,
terpenoids and hybrids thereof -- with roles spanning virulence, chemical ecology, and pharmacology,
from aflatoxins and other mycotoxins to clinically important natural products. The genes responsible
for a given metabolite are typically physically clustered in the genome as a biosynthetic gene
cluster (BGC), and the Minimum Information about a Biosynthetic Gene cluster (MIBiG) standard and
its accompanying database have become the community reference point for recording which BGCs have
been characterised and what they are reported to produce \citep{terlouw2023mibig}.

A MIBiG accession, however, is deliberately a minimal record: cluster boundaries, a product name,
and a small set of controlled fields. The primary evidence behind that record -- which individual
gene was actually disrupted and shown to be required, which metabolite congener was directly
detected versus inferred, under what culture conditions, with what caveats, and whether independent
studies agree -- lives only in the cited literature, in a different format for every paper, and is
not itself queryable. This creates a substantial gap between what is nominally known about a fungal
BGC and what can be programmatically retrieved, compared across BGCs, or checked for internal
disagreement.

\todo{Introduction -- add 2--4 discipline-specific citations here once verified (e.g. antiSMASH/BGC
genome-mining tools, prior fungal-BGC or natural-product knowledge-graph resources, and any of the
609 dossiers' own source papers that motivate specific examples). None are included yet because
fabricating bibliographic details from memory is deliberately avoided; see
\texttt{manuscript/latex/references.bib} for a note on this.}

The FAIR principles -- that data and the tools built on them should be Findable, Accessible,
Interoperable and Reusable \citep{wilkinson2016fair} -- provide a natural design target for closing
this gap: rather than re-reading primary literature for every comparative question, the evidence
itself should be captured once, in a documented schema, with explicit provenance and with
disagreement preserved rather than silently resolved.

Here we present the \textbf{Fungal BGC Atlas}, a resource that applies this design to 609 fungal
BGC dossiers, each linked to a MIBiG accession. Every dossier was built by large-language-model
(LLM)-assisted extraction from the primary literature and MIBiG record, followed by full manual
review of all 609 dossiers by the curator (Methods). The Atlas resolves each dossier into its
constituent genes, metabolites, experiments, measurements, claims and relationships, all linked back
to their supporting sources, and explicitly records the facets, uncertainties and unresolved
normalisation questions that a single-value schema would otherwise discard. We describe the
resource's construction and content, summarise its composition, and present an interactive web
dashboard and archived data/code package that make the Atlas directly usable for cross-BGC
comparison and genome-mining hypothesis generation.

\section{Materials and methods}

\subsection{Data curation}

Each of the 609 dossiers originates from a MIBiG accession and its associated primary reference(s).
Structured fields (BGC identity and boundaries, organism and strain, genes, metabolites, experiments,
measurements, claims, relationships, and source citations) were extracted from the primary literature
and the MIBiG record using Claude (Anthropic) large language models
\todo{Methods -- state the exact Claude model version(s) and extraction date range once confirmed; left as a placeholder rather than guessed}.
Every one of the 609 dossiers was then \textbf{fully manually reviewed} by the curator against the
cited literature before being included in this snapshot; this is a completeness property of the
current release, not a partial-sample audit.

Extraction outputs were passed through a normalisation layer (the \texttt{TERM} vocabulary and
\texttt{NORMALIZATION\_QUEUE} tables) that reconciles free-text labels for experimental conditions,
interventions, assays and taxonomy against a controlled vocabulary where possible. Values that could
not yet be confidently normalised (composite labels, type mismatches, or terms awaiting review) are
retained rather than forced into a canonical bucket, and are tracked explicitly: 882 \texttt{TERM}
rows remained in the normalisation queue at the time of this snapshot (Results). Similarly, where
independent sources disagreed, or a boundary or causal claim was ambiguous, this was recorded as a
\texttt{CONFLICT} row with an explicit issue type and status rather than resolved by editorial
fiat -- 995 such records exist in the current snapshot, of which 220 remain explicitly \emph{open}.

\subsection{Database construction}

All code is released under \texttt{scripts/} in the project repository and is orchestrated by a
single entry point, \texttt{run\_all.py}, so the entire resource is deterministically rebuildable
from the source workbook:

\begin{enumerate}
  \item \textbf{\texttt{parse\_xlsx.py}} reads the 23 data tables and 5 documentation sheets from the
    source workbook and writes one clean CSV per table.
  \item \textbf{\texttt{build\_database.py}} loads the parsed tables into a documented SQLite
    database (\texttt{data/db/fungal\_bgc\_atlas.db}, schema in \texttt{db/schema.sql}), with every
    table's row count asserted against the source workbook's own summary counts as a build-time
    regression check.
  \item \textbf{\texttt{build\_graph.py}} constructs two graphs: a literal knowledge graph of the
    curated subject--predicate--object \texttt{RELATIONSHIP} triples, and an entity-linkage graph
    joining BGC, organism, gene, metabolite and source nodes on \texttt{dossier\_id}, exported as
    GraphML, JSON and CSV edge/node lists.
  \item \textbf{\texttt{compute\_stats.py}} computes every descriptive statistic used in this
    manuscript and in the dashboard directly from the parsed tables, and emits both CSV and
    \LaTeX{}-ready table snippets (\texttt{tables/}).
  \item \textbf{\texttt{make\_figures.py}} renders dated, publication-quality figure snapshots
    (300\,dpi PNG + vector SVG) to \texttt{figures/<date>/}, with an accompanying manifest.
  \item \textbf{\texttt{build\_dashboard.py}} renders the interactive dashboard
    (\texttt{docs/index.html}) described in Results.
\end{enumerate}

The relational schema comprises 23 tables sharing \texttt{dossier\_id} as a common join key (Table~1
of the online Data Dictionary panel; full schema in \texttt{db/schema.sql}), spanning identity and
administration (\texttt{BGC}, \texttt{ORGANISM}, \texttt{STRAIN}), independent categorisation
(\texttt{BGC\_FACET}, 27 axes), molecular content (\texttt{GENE}, \texttt{METABOLITE}), experimental
context (\texttt{EXPERIMENT}, \texttt{TERM}, \texttt{EXPERIMENT\_LINK}), results
(\texttt{MEASUREMENT}, \texttt{OBSERVATION}), interpretive records (\texttt{CLAIM},
\texttt{RELATIONSHIP}, \texttt{CONFLICT}), provenance (\texttt{SOURCE}, \texttt{EVIDENCE\_LINK},
\texttt{ASSERTION}), and curation-audit tables (\texttt{COMPACT\_FACT}, \texttt{COVERAGE},
\texttt{DOSSIER\_SCHEMA}, \texttt{NORMALIZATION\_QUEUE}, \texttt{PROVENANCE\_EXCLUDED},
\texttt{RAW\_ROW}).

The interactive dashboard is a single self-contained HTML page using Plotly for charting and a
small vanilla-JavaScript component for the searchable dossier browser, requiring no server or
build step to view. A GitHub Actions workflow rebuilds the database, graph, statistics, figures and
dashboard on every push to the main branch and deploys the result to GitHub Pages, so the published
resource is always reproducible from the checked-in source data and code.

\section{Results}

\subsection{Resource composition}

The current snapshot (21 August 2026) integrates \textbf{609 fungal BGC dossiers}, each linked to a
MIBiG accession, spanning primary references from 1991 to 2024 (Table~1, Fig.~3). The dossiers cover
\textbf{371 curated organisms} and 1018 distinct strain labels; where taxonomy could be resolved,
organisms span three phyla (167 Ascomycota, 10 Basidiomycota, 1 Mortierellomycota) and 68 genera, and
of the 609 dossiers with a genus-resolved organism (283 of 609; genus is not yet resolved for the
remaining organisms, Discussion), \textit{Aspergillus} accounts for 113 dossiers, \textit{Penicillium}
36, and \textit{Fusarium} 17 (Fig.~1), consistent with the historical emphasis of fungal
natural-product genome mining on these genera. Beyond the cluster level, the Atlas resolves 1572
genes, 807 metabolites, 915 experiments, 936 measurements and 1641 observations (Table~1).

\begin{table}[htbp]
\centering
\caption{Fungal BGC Atlas: resource summary metrics.}
\label{tab:summary}
\begin{tabular}{l l l}
\toprule
Metric & Value & Notes \\
\midrule
BGC dossiers & 609 & All linked to MIBiG accessions \\
Unique organisms & 371 & 3 phyla, 68 genera \\
Strain labels & 1018 &  \\
Genes & 1572 &  \\
Metabolites & 807 &  \\
Experiments & 915 &  \\
Measurements & 936 &  \\
Claims & 1534 &  \\
Relationships (KG triples) & 1264 &  \\
Conflicts / uncertainties & 995 & preserved, not collapsed \\
Sources & 2575 &  \\
Evidence links & 47768 & record <-> source provenance \\
Publication span & 1991-2024 &  \\
\bottomrule
\end{tabular}
\end{table}

\subsection{Biosynthetic and chemical diversity}

MIBiG-reported BGC classes are dominated by polyketide synthase (PKS) and non-ribosomal peptide
synthetase (NRPS) clusters: of the 20 most common raw class labels, ``PKS (Unknown)'' (63 dossiers),
``NRPS (Type I)'' (57) and ``PKS (Type I)'' (38) are the largest single categories, with a further 19
dossiers combining NRPS and PKS machinery and 18 assigned to terpene biosynthesis (Table~3, Fig.~2).
Beyond this single class label, each BGC is independently characterised along 27 facet axes spanning
biosynthesis, chemistry, ecology, genome context and activation conditions (\texttt{BGC\_FACET}, 7800
rows), allowing a cluster to be queried simultaneously by, for example, chemical class, host
organism ecology, and the developmental state under which it is known to be expressed -- without
collapsing these into a single forced category.

\subsection{Knowledge graph of curated relationships}

The \texttt{RELATIONSHIP} table records 1264 curated statements of the form
subject--predicate--object (e.g.\ gene \emph{encoded\_by} protein, cluster \emph{produces}
metabolite, intervention \emph{required\_for} phenotype); 994 of these have a complete triple and
are rendered as the Atlas knowledge graph (1870 entity nodes, 994 typed edges), which can be
filtered interactively by predicate in the dashboard (Fig.~4 of the online dashboard; Materials and
methods). The most common predicates are \emph{produces} (60 instances), \emph{required\_for} (31),
\emph{synthesizes} (28), \emph{cyclizes} (27) and \emph{encoded\_by} (26), reflecting the underlying
emphasis on causal, pathway-level evidence rather than purely descriptive annotation.

\subsection{Evidence, provenance and preserved uncertainty}

Every claim, measurement, gene function and relationship in the Atlas is traceable to a
\texttt{SOURCE} record (2575 sources) through an explicit \texttt{EVIDENCE\_LINK} (47\,768 links),
so a user can move from any field back to the literature or database record that supports it
(Fig.~4). Rather than resolving disagreement editorially, 995 \texttt{CONFLICT} records preserve
boundary disputes, strain- or assay-dependent discordance, and unresolved causal questions as
first-class, queryable data; the most common issue types are ecological-role ambiguity (20 records),
BGC boundary-scope disagreement (18) and depth-of-evidence questions about gene function (13)
(Table~4, Fig.~5). Of the 995 conflicts, 220 remain explicitly open and 199 are marked
resolved-for-this-Atlas, with the remainder spanning provenance- or curation-specific resolutions --
the full distribution is available in \texttt{supplementary/supp\_table9\_conflict\_status.csv}.

\subsection{Interactive dashboard}

The Atlas is published as a single self-contained web page
(\url{https://dr-richard-barker.github.io/fungal-bgc-atlas/}) with nine linked panels: an
\textbf{Overview} of headline metrics and per-dossier evidence-coverage completeness (Fig.~6); a
\textbf{Taxonomy Explorer} (phylum--class--order--genus composition); a \textbf{Chemistry \&
Biosynthesis} panel switching between BGC class, chemical class, product family and biological role;
a \textbf{Facet Explorer} over all 27 facet axes; the \textbf{Knowledge Graph} described above; an
\textbf{Evidence \& Provenance} panel (source types, publication timeline, evidence-link density); a
\textbf{Conflicts \& Uncertainty} panel; a searchable \textbf{Dossier Browser} that assembles, for any
of the 609 BGCs, its full cross-table record (organism, strain, genes, metabolites, experiments,
measurements, claims, relationships, conflicts and sources, joined on \texttt{dossier\_id}); and a
live \textbf{Data Dictionary} rendered directly from the source workbook's own schema documentation.

\section{Discussion}

The Fungal BGC Atlas is designed to complement, not replace, MIBiG: MIBiG remains the authoritative,
minimal, sequence-anchored record of a BGC, while the Atlas captures the surrounding evidence --
which gene was shown to matter, which metabolite was actually detected, under what conditions, by
whom, and with what caveats -- as structured, queryable, and independently citable data. By
resolving each dossier into genes, metabolites, experiments, measurements, claims and typed
relationships, all linked back to source records, cross-cluster questions that previously required
re-reading the primary literature (which predicates most often co-occur with a given chemical class;
which BGCs have unresolved boundary disputes; which experimental conditions are associated with a
given activation state) become directly queryable against the SQLite database or the knowledge
graph.

Three design choices follow the FAIR principles \citep{wilkinson2016fair} directly. First,
\textbf{Interoperability} is served by a documented relational schema in which every table shares
\texttt{dossier\_id} as a join key, so the entity-linkage and knowledge-graph views in the dashboard
are direct database queries rather than bespoke integration code. Second, \textbf{Reusability} is
protected by preserving disagreement and unresolved normalisation explicitly (995 conflict records,
882 open normalisation-queue rows) rather than silently forcing a single value, so a downstream user
can see exactly where the curated record is contested or provisional. Third, \textbf{Findability and
Accessibility} are served by packaging the database, knowledge graph, figures, tables, supplementary
material and code together for Zenodo archival under an open licence, with a live interactive
dashboard as the primary point of entry.

\subsection{Limitations}

Two limitations follow directly from the current snapshot's construction and are stated explicitly
rather than hidden. First, taxonomic resolution is incomplete: phylum could be resolved for 178 of
371 organisms and genus for organisms underlying 283 of 609 dossiers at this snapshot; the remainder
await taxonomy normalisation and are tracked in \texttt{NORMALIZATION\_QUEUE} rather than assigned a
placeholder value. Second, while every dossier was fully manually reviewed by the curator, this
release reflects a single-curator review pass; broadening review to multiple independent curators,
and formalising inter-curator agreement on the \texttt{CONFLICT} and normalisation-status fields,
is a natural next step and is discussed as future work below.
\todo{Discussion -- add the exact LLM model version(s)/dates here alongside the limitation statement once confirmed (Methods).}

\subsection{Future work}

The build pipeline (Materials and methods) is designed so that on-boarding additional BGCs is an
append-and-rebuild operation rather than a schema change: new dossiers extend the same 23-table
schema, are validated by the same row-count and referential-integrity checks, and immediately appear
in the dashboard and knowledge graph. Natural extensions include: closing the normalisation queue
(882 open \texttt{TERM} rows) through continued curator review; resolving the remaining open
conflicts (220 of 995) as new evidence becomes available; cross-linking dossiers to external
resources beyond MIBiG (e.g.\ chemical structure databases for the 807 curated metabolites); and
opening the curation workflow to community contribution with the same evidence-linked,
conflict-preserving schema enforced at ingestion.

\section{Data availability}

The Fungal BGC Atlas database, knowledge graph, source code, figures, tables and supplementary
material are freely available under the Creative Commons Attribution 4.0 International licence
(CC-BY-4.0) at:

\begin{itemize}
  \item \textbf{Interactive dashboard:} \url{https://dr-richard-barker.github.io/fungal-bgc-atlas/}
  \item \textbf{Source repository (code + data + manuscript):}
    \url{https://github.com/dr-richard-barker/fungal-bgc-atlas}
  \item \textbf{Archived release (Zenodo):} \todo{Zenodo DOI -- to be assigned when the repository is
    deposited; this manuscript is prepared alongside the Zenodo packaging (\texttt{.zenodo.json},
    \texttt{CITATION.cff}) but the deposit itself is a separate step.}
\end{itemize}

All code used to parse the source workbook, build the relational database and knowledge graph,
compute the statistics reported here, render the figures, and build the dashboard is provided in
\texttt{scripts/} and orchestrated by \texttt{run\_all.py}, so the entire resource is
deterministically reproducible from the checked-in source data.

\section*{Availability of source code and requirements}
\begin{itemize}
  \item Project name: Fungal BGC Atlas
  \item Project home page: \url{https://github.com/dr-richard-barker/fungal-bgc-atlas}
  \item Operating system(s): Platform independent
  \item Programming language: Python 3.9+ (build pipeline); HTML/JavaScript (dashboard)
  \item License: CC-BY-4.0
\end{itemize}


\section*{Conflict of interest}
None declared.

\section*{Funding}
\todo{Funding -- confirm before submission (e.g. institutional support, grant numbers, or ``no external funding was received for this work'').}

\section*{Acknowledgements}
\todo{Acknowledgements -- confirm before submission (e.g. curators, reviewers, computing resources).}

\bibliographystyle{unsrtnat}
\bibliography{references}

\end{document}
